\section{Hardware Software Co-design}


Because industrial hardware vendors' platforms are not open source, security proposals building on top of their hardware suffer from several limitations as they are limited to software development.
First, it is virtually impossible to make strong security arguments without a clear (open-source) view of the hardware itself.
It is also difficult to repurpose existing mechanisms for security as they can lead to small but important mistakes.
These elements limit security and constrains the threat model of these proposals as it is inherently difficult to protect against novel attacks that exploit microarchitecture without sound hardware support.

On the other hand, the only previous work \cite{bourgeat2019mi6} to address our threat model implemented all its defense mechanisms in hardware and enabled very few interactions with the firmware or software layer.
This hardware-only design provides very coarse-grained mechanisms for which security is easy to reason about but suffers from limited performance.
For instance, flushing is systematically triggered by the hardware for every single trap handling and every partitioning scheme is statically defined in hardware.
It also aimed at evaluating parts of the proposed design, only modeling the performance cost of some secure mechanisms and did not build a real platform.

Given these insights, we believe hardware software co-design is required to build secure and usable enclaves. 
We present two other mechanisms on top of secure shared memory that aim at refining MI6 coarse-grained hardware mechanisms to improve performance and enclave usability while maintaining its security guarantees.
We analyze the security of these new mechanisms in Section \ref{sec:secannalysis} and evaluate their performance implications in Section \ref{sec:eval}.



\subsection{Dynamic Cache Partitioning}
\label{sec:dynamiccache}
We aim at demonstrating a secure enclave architecture for a different class of machines and constrain ourselves to building a 2-core prototype with a 1MB LLC.
This conveniently makes it possible to run our entire design on an FPGA.
Unfortunately, this makes static partitioning (as proposed in MI6) impractical as each of the 64 DRAM regions would be allocated a 16KB cache slice: a slice even smaller than the 64KB necessary for the inclusion of the cores' private L1 caches (D1 and I1).
This dramatically impacts performance.
Moreover, using way partitioning \cite{kiriansky2018dawg} is not possible as we need to instantiate 64 regions in the cache but the LLC has only 16 ways, and Ironhide's \cite{ironhide} strategy assumes limited reconfigurations.
We propose a new, simple, dynamic set partitioning scheme that only requires minimal changes in the hardware. 
We start with a high-level description of the idea, and then we detail the family of set-index functions we use.

The goal is to let the OS request the SM to configure the cache to allocate a variable amount of LLC for each DRAM region. Our LLC contains 1024 sets, and we propose that each DRAM region may be given a number of sets equal to any power of two up to 512 sets.
After controlling the validity of the partitioning (i.e., that it fits in the LLC), the SM conveys the new mapping to the hardware by writing the sizes of each LLC slice in specialized CSRs.
As $512=2^9$, and $9$ can be represented with $4$ bits,  we only need 4 bits per DRAM region to describe the size of LLC that we allocate for the DRAM region. This gives a total of $64 \times 4$ bits. We can expose these configuration registers through 4 new CSRs.
We cannot change these configurations at arbitrary times, as it would change the set index mapping function for in-flight requests, breaking the cache.
Intuitively, changing the set-index mapping function should only be done whenever the cache is fully flushed and no memory requests are in flight.
To implement such a switching policy, we suggest using 9 (= 4 + 4 + 1) CSRs, and use a form of \emph{double-buffering}.
The two sets of CSRs would hold both the configuration currently in-use and the future configuration. 
The SM only modifies the future configuration and never touches the configuration currently in-use. 
When the SM finishes modifying the 4 future configuration CSRs, it can set the 9th CSR to toggle which of the two sets of 4 CSRs is active to describe the set-index mapping function.
Because the LLC is shared (there is a coherency protocol), the toggling is itself non-trivial.
When the toggling is requested the processor must in sequence: 1. block any new memory requests from entering the memory subsystem, 2. drain all outstanding memory requests, 3. flush the L1s and the LLC (as the caches are inclusive), and 4. toggle the CSRs in use to describe the set-index mapping function.
Fortunately, this costly operation is only required when changing the LLC layout which mostly happens on enclave creation or when launching an untrusted application with an important memory footprint.

Finally, we describe the family of set-index mapping functions used to implement our dynamic set partitioning scheme.
Remember that our scheme aims at splitting the 1024 set indexes of the cache among the 64 different DRAM regions and making sure two addresses from different regions are never assigned to the same index range, i.e., LLC slice.
We represent the algorithm used to compute the set index for a given address in Figure \ref{fig:dynamicpartitioning}.
The idea is to store the sizes of the 64 index ranges in a 256 + 1024 bits register called the set-index range table, along with the base of each range.
The bases are computed during the partition switch and are equal to the sum of the size of previous ranges.
When the LLC receives a request, it will compute the set index for the corresponding address.
It will first extract a region ID from the address upper bits and obtain the corresponding index range size and bound by looking up the set index range table \circled{1}.
It will then read the address original set index, located in the middle of the address, set the upper bits to zero depending on the range size \circled{2} and add the base index \circled{3}.
This effectively restricts the new set index to the corresponding index range, partitioning the LLC \circled{4}.
Note that for this to work, we must extend the length of the tags used by the LLC to include the entire original index to cover the case where an index range is of size 1 and all addresses from the same region map to the same set.

We show the significant advantages of this dynamic LLC partitioning in Section \ref{sec:eval}.


\begin{figure}[t]
\centering
\includegraphics[width=\columnwidth]{figures/partitioning.png}
\caption{LLC set indexing scheme. The function takes an address as an input and output a set index in the LLC slice allocated to its memory region.}
%\vspace{8pt}
\label{fig:dynamicpartitioning}
\end{figure}


\subsection{Relaxing the Flushing Policy}
Because we build the entire software infrastructure required to run enclaves, we observe that some security domains considered by MI6, especially the enclave and some parts of the SM, do not have secrets to keep from each other (see our security analysis in Section \ref{sec:secannalysis} for more details).
As a result, we disable flushing on these context switches.
Enclaves can now access the SM without paying the prohibitive cost of a microarchitectural flush.
This opens the door to the use of frequent software or timer interrupts.

To disable flushing selectively, we modify the hardware to deactivate systematic flushing on traps and SM-return instructions.
We add a CSR that, when accessed, triggers the microarchitectural flushing.
We modify the SM to trigger the flushing mechanism on required context switches but not on most traps and SM-calls from the enclave.
We evaluate the performance implication of our mechanism in Section \ref{sec:evalhsmecha}. 


\subsection{Partitioning MSHRs and the LLC arbiter}

MI6 modeled the performance impact of partitioning of the MSHRs and the arbiter located in the LLC.
We implement both security mechanisms.
We allocate half of the MSHRs to the first core and the second half to the second. Our arbiter then serves requests from the two different MSHR sets using a fair round-robin arbiter to avoid contention between the two security domains.
Because these two mechanisms only marginally affect the latency of the LLC for a 2-core design, we observe negligible ($\approx$ 0.1\%) performance impact.